% Part:sets-functions-relations
% Chapter: size-of-sets
% Section: reduction-alt

\documentclass[../../../include/open-logic-section]{subfiles}

\begin{document}

\olfileid{sfr}{siz}{red-alt}

\olsection{راکمونه}

\begin{editorial}
  دا برخه په راکمونه د شمېر نۀ وړ والے ثابتوي، او پايلې ئې د
  \olref[nen-alt]{sec} له پايلو سره سمون خوري۔ يوه بله، لږه
  تفصيلي بڼه چې د \olref[nen]{sec} له پايلو سره سمون خوري، په
  \olref[red]{sec} کښې ورکړل شوې ده۔
\end{editorial}

موږ په يوه قطري استدلال وښودل چې $\Bin^\omega$ !!{nonenumerable}
دے۔ په ورته قطري استدلال مو وښودل چې $\Pow{\Nat}$
!!{nonenumerable} دے۔ خو د $\Pow{\Nat}$ د !!{nonenumerable}
کېدو يوه بله لار دا ده: وښيئ چې \emph{که $\Pow{\Nat}$
!!{enumerable} وي، نو $\Bin^\omega$ هم !!{enumerable} دے}۔ ځکه
چې پوهېږو $\Bin^\omega$ !!{nonenumerable} دے، دا به وښيي چې
$\Pow{\Nat}$ هم همداسې دے۔

دې ته د يوې مسئلې بلې ته \emph{راکمونه} ويل کېږي۔ په دې حالت
کښې د $\Bin^\omega$ د شمېرلو مسئله د $\Pow{\Nat}$ د شمېرلو
مسئلې ته راکموو۔ د وروستۍ مسئلې يو حل، يعنې د $\Pow{\Nat}$ يوه
شمېرنه، به د لومړۍ مسئلې حل، يعنې د $\Bin^\omega$ يوه شمېرنه،
راکړي۔

د يوه سټ~$B$ د شمېرلو مسئله د يوه سټ~$A$ د شمېرلو مسئلې ته د
راکمولو دپاره، د~$A$ د يوې شمېرنې د~$B$ په شمېرنه د اړولو يوه
لار ورکوو۔ د دې تر ټولو اسانه لار د يوې !!a{surjection} تابع
$f\colon A \to B$ تعريفول دي۔ که $x_1$، $x_2$، \dots{} د~$A$
شمېرنه وي، نو $f(x_1)$، $f(x_2)$، \dots{} به د~$B$ شمېرنه وي۔
زموږ په حالت کښې د يوې !!a{surjection} تابع
$f\colon \Pow{\Nat} \to \Bin^\omega$ په لټه کښې يو۔

\begin{prob}
وښيئ چې که يوه !!{injective} تابع $g\colon B \to A$ وي او
$B$~!!{nonenumerable} وي، نو $A$ هم همداسې دے۔ دا په ښودلو
وکړئ چې څنګه د~$g$ په کارولو د~$A$ يوه شمېرنه د~$B$ په شمېرنه
اړولے شئ۔
\end{prob}

\begin{proof}[په راکمونه د {\olref[nen-alt]{thm:nonenum-pownat}} ثبوت]
فرض کړئ چې $\Pow{\Nat}$ !!{enumerable} وے، او له دې امله د هغه
يوه شمېرنه $N_{1}$، $N_{2}$، $N_{3}$، \dots وے۔

تابع $f \colon \Pow{\Nat} \to \Bin^\omega$ داسې تعريف کړئ چې
$f(N)$ هغه تسلسل $s$ وي چې $s(n) = 1$ هله او يوازې هله چې
$n \in N$ وي، او په بل حالت کښې $s(n) = 0$ وي۔
\textbf{د ژباړې د نسخې سمون (OLSIZ-010):} په انګرېزي نثر کښې
د تسلسل بې تړاوه شاخص د خپل سري ورودي سټ له نښې سره بدل شو، څو
د تابع تعريف روښانه وي۔

دا په څرګند ډول تابع تعريفوي، ځکه کله چې $N \subseteq \Nat$ وي،
هر $n \in \Nat$ يا د~$N$ يو !!a{element} وي يا نۀ وي۔ د بېلګې
په ډول، د طبيعي جفتو عددونو سټ
$2\Nat = \Setabs{2n}{n \in \Nat} = \{0,2, 4, 6, \dots\}$
تسلسل $1010101\dots$ ته نقشېږي، $\emptyset$ سټ $0000\dots$ ته او پخپله
سټ $\Nat$، $1111\dots$ ته نقشېږي۔

دا !!{surjective} هم ده: د $0$ګانو او $1$ګانو هر تسلسل د طبيعي
عددونو له يوه سټ سره سمون خوري، يعنې له هغه سټ سره چې غړي ئې
هغه طبيعي عددونه دي کوم چې د تسلسل د~$1$ لرونکو مقامونو سره
سمون خوري۔ په لا دقيقه توګه، که $s \in \Bin^\omega$ وي، نو
$N \subseteq \Nat$ داسې تعريف کړئ:
\[
N = \Setabs{n \in \Nat}{s(n) = 1}
\]
بيا $f(N) = s$، لکه چې د~$f$ تعريف ته په کتلو تاييدولے شو۔

اوس دا لړ په پام کښې ونيسئ:
\[
f(N_1), f(N_2), f(N_3), \dots
\]
ځکه چې $f$ !!{surjective} ده، د $\Bin^\omega$ هر غړے بايد د کوم
ورودي دپاره د~$f$ د قيمت په توګه راشي، او نو بايد په لړ کښې راشي۔
له دې امله دا لړ بايد ټول~$\Bin^\omega$ وشمېري۔

نو که $\Pow{\Nat}$ !!{enumerable} وے، $\Bin^\omega$ به
!!{enumerable} وے۔ خو $\Bin^\omega$ !!{nonenumerable} دے
(\olref[nen-alt]{thm:nonenum-bin-omega})۔ له دې امله
$\Pow{\Nat}$ !!{nonenumerable} دے۔
\end{proof}

%\begin{explain}
%د راکمونې په لوري کښې ګډوډېدل اسانه دي۔ د بېلګې په ډول، يوه
%!!a{surjective} تابع $g \colon \Bin^\omega \to X$ دا \emph{نۀ}
%ثابتوي چې $X$ !!{nonenumerable} دے۔ (تابع $g
%\colon \Bin^\omega \to \Bin$ په پام کښې ونيسئ چې په
%$g(s) = s(1)$ تعريف شوې؛ دا تابع د $0$ګانو او $1$ګانو يو تسلسل
%د هغه لومړي !!{element} ته نقشوي۔ دا پر هدف سټ بشپړه ده، ځکه ځينې
%تسلسلونه په $0$ او ځينې په $1$ پيلېږي۔ خو $\Bin$ متناهي دے۔)
%دا هم په ياد ولرئ چې تابع~$f$ بايد پر هدف سټ بشپړه وي، ګنې
%استدلال مخکښې نۀ ځي: بيا به دا ډاډ نۀ وي چې $f(x_1)$، $f(x_2)$،
%\dots{} د~$Y$ ټول !!{element}s لري۔ د بېلګې په ډول، تابع
%$h\colon \Nat \to
%\Bin^\omega$ په لاندې ډول تعريف کړئ:
%\[
%h(n) = \underbrace{000\dots0}_{\text{$n$ $0$ګانې}}
%\]
%دا يوه تابع ده، خو $\Nat$ !!{enumerable} دے۔
%\end{explain}

\begin{prob}\label{sfr:siz:red:prob:nat-nat}
  په راکمونې استدلال وښيئ چې د ټولو تابعو $f\colon \Nat \to \Nat$
  سټ~$X$ !!{nonenumerable} دے۔ (لارښوونه: له $X$ څخه
  ~$\Bin^\omega$ ته يوه پر هدف سټ بشپړه تابع ورکړئ۔)
\end{prob}

\begin{prob}
په راکمونې استدلال وښيئ چې د طبيعي عددونو د جوړو د
\emph{سټونو} سټ، يعنې $\Pow{\Nat \times \Nat}$،
!!{nonenumerable} دے۔
\end{prob}

\begin{prob}
په راکمونې استدلال وښيئ چې $\Nat^\omega$، يعنې د طبيعي عددونو
د نامتناهي تسلسلونو سټ، !!{nonenumerable} دے۔
\end{prob}

%\begin{prob}
%فرض کړئ $P$ له $\Nat$ څخه سټ~$\{0\}$ ته د تابعو سټ دے، او $Q$
%له مثبتو صحيح عددونو څخه سټ~$\{0\}$ ته د \emph{جزوي} تابعو سټ
%دے۔ وښيئ چې $P$~!!{enumerable} دے او $Q$ نۀ دے۔ (لارښوونه: د
%$\Bin^\omega$ د شمېرلو مسئله د~$Q$ شمېرلو ته راکمه کړئ۔)
%\end{prob}

\begin{prob}
فرض کړئ $S$ له $\Nat$ څخه سټ~$\{0,1\}$ ته د ټولو
!!{surjection}s سټ دے؛ يعنې $S$ له ټولو !!{surjection}s
~$f \colon \Nat \to \Bin$ څخه جوړ دے۔ وښيئ چې $S$
!!{nonenumerable} دے۔
\end{prob}

\begin{prob}
وښيئ چې د ټولو حقيقي عددونو سټ~$\Real$ !!{nonenumerable} دے۔
\end{prob}

\end{document}
